\chapter{项目背景与行业痛点}

\section{农业数字化转型的行业背景}

在乡村振兴战略与数字乡村建设的双重推动下，我国农业正经历从「传统经验农业」向「数据智能农业」的深刻转型。以 5G、物联网、人工智能、遥感为代表的新一代信息技术加速渗透农业生产全链条，智慧农业、精准农业成为政策与资本共同聚焦的方向。病虫害防治作为直接影响粮食安全与农民收益的核心环节，其数字化、智能化水平却长期滞后——大量小农户仍停留在「凭经验、问熟人、看天吃饭」的阶段。这一「最后一公里」的技术鸿沟，正是本项目切入的窗口。

\section{小农户病虫害防治的真实痛点}

\subsection{经验依赖导致识别不准、用药错误}

小农户遇到病虫害，第一反应是「凭经验」或「问熟人」。然而病虫害种类繁多、症状相似（如稻瘟病与胡麻叶斑病早期病斑高度相近），仅凭肉眼与经验极易误判。识别不准直接导致\hlgold{「用错药、用重药、用贵药」}：既造成农资浪费与药害残留，又延误正确防治，甚至引发食品安全与环境污染问题。

\subsection{防治窗口期短，易错过最佳时机}

绝大多数病虫害的防治存在明确的「\hlgold{窗口期}」——错过几天，防效大幅下降，轻则减产、重则绝收。但农户往往缺乏对「本区域、本季节、本作物高发病虫害」的预判能力，只能被动应对，等到肉眼可见的严重症状出现时，往往已错过最佳防治时机。

\begin{figure}[H]
  \centering
  \includegraphics[width=0.86\linewidth]{painpoints.png}
  \caption{小农户病虫害防治核心痛点（调研占比）}
\end{figure}

\section{现有解决方案的局限性}

\subsection{传统农技服务：覆盖不足、响应滞后}

基层农技推广体系长期面临「人少、面广、任务重」的困境，一名农技人员往往要服务数千户农户，难以做到及时、精准、一对一的诊断指导。农资店的「土专家」虽有一定经验，但水平参差、利益驱动下常有推销倾向，难以提供客观中立的防治方案。

\subsection{现有农业APP：功能冗余、门槛过高}

市面上的农业综合类 APP 普遍追求「大而全」，功能堆叠、界面复杂、专业术语密集，与真实的小农户使用习惯严重脱节。数据显示，农户智能手机使用率已达 76\%，但真正将其用于农业生产管理的极少，\hlgold{约三分之二的受访者从未关注过农业类公众号}。这充分说明：不是农户「不用手机」，而是现有产品「门槛太高、不够好用」。

\section{项目发起的必要性与社会价值}

综上所述，小农户病虫害防治面临「\hlred{识别难、方案缺、预警弱、门槛高}」四大结构性矛盾，而传统农技服务与现有农业 APP 均未能有效破解。本项目以「极简、轻量、落地」为核心原则，用一部手机、一张照片，把专业植保能力送到亿万小农户手中，既是填补市场空白的商业机会，更是\hlgreen{降低农业生产风险、保障粮食安全、助力农民增收}的社会价值工程。
